\documentclass{beamer}
\usetheme{Madrid}
\usecolortheme{beaver}

\usepackage[utf8]{inputenc}
\usepackage[T1]{fontenc}
\usepackage[portuguese]{babel}
\usepackage{amsmath,amssymb,amsthm}
\usepackage{microtype}

\title[Teorema da Urna (Ballot)]{Exemplo: Teorema da Urna (Ballot Theorem)}

\author[MCBM022-23]{MCBM022-23 Introdução aos Processos Estocásticos }

\date{UFABC 2025-3}


\begin{document}

\begin{frame}
  \titlepage
\end{frame}


\begin{frame}{Enunciado do problema}
  \begin{block}{Problema}
    Dois candidatos concorrem numa eleição. O candidato A obtém $a$ votos e o candidato B obtém $b$ votos, com $a>b$. Os $n=a+b$ votos são contados numa ordem aleatória (permutação uniforme).  
    Qual é a probabilidade de que A esteja sempre à frente ao longo da contagem?
  \end{block}
  
\end{frame}

\begin{frame}{Definições e notação}
  \begin{itemize}\itemsep=1cm
  \item $n=a+b$ total de votos.
  \item $S_k$ = \( (\#\text{votos de A} - \#\text{votos de B} )\) após $k$
    votos contados

    (logo $S_n=a-b$). Note que $S_k$ pode ser negativo.
  \item Para $0\le k\le n-1$ defina
    \[
      X_k=\frac{S_{n-k}}{\,n-k\,}.
    \]
  \end{itemize}
  
  
  \textbf{$(X_k)_k$ é martingal em relação a
    $(\sigma(X_0,\dots,X_{k-1}))_k$.}
  
\end{frame}
% -----------------------------------------------
\begin{frame}{Por que $(X_k)$ é  martingal?}
  Queremos mostrar: \(\mathbb{E}[X_k\mid X_0,\dots,X_{k-1}]=X_{k-1}.\)
  
  \medskip Condicionar em $X_0,\dots,X_{k-1}$ equivale a condicionar em
  \(S_n,S_{n-1},\dots,S_{n-k+1}\).
  
  \medskip
  Dado $S_{n-k+1}$, com $n+k-1$ votos contados
  \begin{itemize}
  \item A teve
    \(\dfrac{n-k+1+S_{n-k+1}}{2}\) votos e 
  \item B teve
    \(\dfrac{n-k+1-S_{n-k+1}}{2}\) votos.
  \end{itemize}

  \medskip

  
  $$
  S_{n-k} =
  \begin{cases}
    S_{n-k+1} +1 & \text{ se o voto foi  em }B\\
    S_{n-k+1} -1 & \text{ se o voto foi em }A
  \end{cases}
  $$
  
  
\end{frame}
%-----------------------------------------------
\begin{frame}{Por que $(X_k)$ é  martingal?}

  O voto foi  escolhido uniformemente 
  \[
    \begin{split}
    \mathbb{E}[S_{n-k}\mid S_{n-k+1}]
    &=
    (\,S_{n-k+1}+1\,) \cdot \frac{n-k+1-S_{n-k+1}}{2(n-k+1)}
    + \\ &
    (\,S_{n-k+1}-1\,)\cdot \frac{n-k+1+S_{n-k+1}}{2(n-k+1)}\\       
    = \frac{n-k}{n-k+1}\,S_{n-k+1}&.
  \end{split}
  \]

  Assim
  \[
    \mathbb{E}\!\left[\frac{S_{n-k}}{n-k}\Bigm\vert S_n,\dots,S_{n-k+1}\right]
    =\frac{S_{n-k+1}}{n-k+1},
  \]
  ou seja \(\mathbb{E}[X_k\mid X_0,\dots,X_{k-1}]=X_{k-1}.\)
\end{frame}

\begin{frame}{Aplicação do teorema de paragem}
  \begin{itemize}\itemsep=1cm
  \item Defina o tempo de parada \(T\) como o menor índice \(k\) tal
    que \(X_k=0\), se existir; caso contrário \(T=n-1\).
    \item Observações:
      \begin{itemize}\itemsep=.5cm
        \item $T$ é um tempo de parada limitado (\(T\le n-1\)), então as hipóteses do teorema de paragem (para martingales limitadas) se aplicam.
        \item Pelo teorema de paragem temos \(\mathbb{E}[X_T]=\mathbb{E}[X_0]=\dfrac{a-b}{a+b}.\)
      \end{itemize}
  \end{itemize}
\end{frame}

\begin{frame}{Análise dos casos}
  Existem duas situações mutuamente exclusivas:

  \medskip
  \textbf{Caso 1: A lidera durante toda a contagem.}
  \begin{itemize}
    \item Então todos os \(S_{n-k}>0\) para \(0\le k\le n-1\), logo todos \(X_k>0\).
    \item Nesse caso \(T=n-1\) e \(X_T=X_{n-1}=S_1/1=1\) (A recebe o primeiro voto).
  \end{itemize}

  \medskip
  \textbf{Caso 2: A não lidera durante toda a contagem.}
  \begin{itemize}
    \item Então existe algum índice em que \(S_k\le 0\). Pelo raciocínio, haverá algum \(k<n-1\) com \(X_k=0\), então \(T< n-1\) e \(X_T=0\).
  \end{itemize}
\end{frame}

\begin{frame}{Concluir a probabilidade}
  Como \(X_T\) só pode assumir valores 0 ou 1 nos cenários acima,
  \[
    \mathbb{E}[X_T]=1\cdot\Pr(\text{Caso 1}) + 0\cdot\Pr(\text{Caso 2}) = \Pr(\text{A sempre à frente}).
  \]
  Mas já vimos que \(\mathbb{E}[X_T]=\mathbb{E}[X_0]=\dfrac{a-b}{a+b}\). Portanto
  \[
    \Pr(\text{A sempre à frente})=\frac{a-b}{a+b}.
  \]
  Isso conclui a prova via martingal.
\end{frame}


\end{document}
